\subsection{معماری و پیاده‌سازی}
پیاده‌سازی ماژول در فایل \lr{Codes/framework/compression\_engine/come\_compressor.py} مستقر بوده و شامل اجزای مهندسی زیر است:
\begin{itemize}
    \item \textbf{کلاس \lr{CoMeBridge}:} پل پروجکشن خطی که به صورت چندریختی انواع خروجی‌های پایتورچ شامل \lr{Tensor}، \lr{Tuple}، \lr{Dict} و \lr{BaseModelOutput} را مدیریت کرده و عملگر درون‌یابی را منحصراً روی تانسور حالت‌های پنهان اعمال می‌نماید.
    \item \textbf{کلاس \lr{CoMeBridgedLayer}:} پوسته‌ای ساختاری که لایه ترنسفورمر نگه داشته شده را دربرگرفته و بلافاصله پس از اجرای گذر رو به جلو، پل درون‌یابی را فراخوانی می‌کند.
    \item \textbf{کلاس \lr{CoMeLoRALinear}:} آداپتورهای کم‌مرتبه \lr{LoRA} کمکی تعبیه‌شده در لایه‌های نگه داشته شده تا فرآیند ریکاوری سریع طی چند گام بهینه‌سازی را تسهیل نمایند.
    \item \textbf{موتور فشرده‌ساز \lr{CoMeCompressor}:} مدیریت کل چرخه فشرده‌سازی، کالیبراسیون مجدد لایه‌های نرمال‌سازی (\lr{recalibrate\_layernorms}) و اتصال پل‌ها به ماژول‌های \lr{nn.ModuleList}.
\end{itemize}
